\documentclass[journal,onecolumn]{IEEEtran}

\usepackage{cite}
\usepackage{amsmath}

\begin{document}

\title{EEC289Q Project Proposal: Efficient Power Analysis Resilience for Neural Networks}

\author{Luna~Park,
        Anjani~Gadkari
}

% make the title area
\maketitle
\section{Problem and Motivation}
Deep neural network (DNN) models are increasingly deployed into security-critical domains including fraud detection and malware 
identification~\cite{alarfajCreditCardFraud2022, vinayakumarRobustIntelligentMalware2019}. This motivates adversaries to locate DNN 
vulnerabilities through model extraction and weight stealing~\cite{chakrabortySurveyAdversarialAttacks2021,rakinDeepStealAdvancedModel2022}. Recent literature has proven that DNN hardware is vulernable to power 
analysis side channel attacks(SCA), which allows adversaries to perform weight stealing on such 
models~\cite{majiLeakyNetsRecovering2021}. While there is significant research in the domain of power analysis resilient hardware, 
established implementations are generally dedicated to protecting cryptographic 
operations~\cite{heHowSecureYour2017,zhaoSideChannelSecurity2023,srivastavaSCARPowerSideChannel2024a}. These algorithms, such as 
AES, require precise comuputation to follow well-established standardizations, which then incurs significant overhead when protecting. 
In contrast, machine learning models are robust to errors in data and computation. This allows approximate computing paradigms to make 
significant improvements in hardware size, performance, and energy usage with minimal loss in accuracy. 

One of the most effective techniques which takes advantage of this tolerance is quantization. Quantization for DNN's reduces 
the bit size of weights and activations, thereby reducing the precision but increasing efficiency when deployed~\cite{rokhComprehensiveSurveyModel2023}. Binarized
Neural Networks (BNNs) are the most efficient form of quantization, constraining weights and activations to a single bit such
that $w, a \in \{-1, 1\}$~\cite{courbariauxBinarizedNeuralNetworks2016}. The single bit nature of BNNs then enables new hardware designs for inference—for example, 
BNN multiplication can be performed with a single XNOR. In this project, we want to exploit the mathematical properties of BNNs
to build power SCA resistant neural network hardware which requires minimal overhead.


\section{Design Choices}
Our design focuses on prioritizing a lightweight and secure architecture for accelerators by utilizing BNNs. 
By using single-bit weights and activations, BNNs significantly reduce memory requirements and computation complexity, 
and can be implemented on resource-constrained hardware. To ensure resistance against SCA attacks, we will implement 
gate-level masking during MAC operations. In this method, every sensitive bit is split into two or more shares such that 
$$x = x_1 \oplus x_2,$$
where $x_1$ is a randomly generated bit. In BNNs, MAC operations can be performed solely through XNOR and AND gates. 
Because the XNOR operation is linear, shares can be processed independently of each other~\cite{trichinaSmallSizeLow2005}. For two secret bits $a$ and $b$
and the operation $z = a \odot b$, the value can be calculated as follows: 
$$z=(a_1 \oplus a_2)\odot(b_1 \oplus b_2) = \overline{(a_1 \oplus b_1)} \oplus (a_2 \oplus b_2).$$
However, AND gates are non-linear, and processing the operation results in shares appearing across multiple terms:
$$z=(a_1 \oplus a_2)\cdot(b_1 \oplus b_2) = (a_1 \cdot b_1) \oplus (a_1 \cdot b_2) \oplus (a_2 \cdot b_1) \oplus (a_2 \cdot b_2)$$

Terms that contain both the secret share bits and masks such as $a_1 \cdot b_2$ are prone to leaking the value of the secret bit. 
To mitigate this, we replace the standard AND-gates with a glitch-resistant implementation of Trichina AND gates~\cite{dubeyBoMaNetBooleanMasking2020}. This implementation 
processes shares with additional random bits to ensure that intermediate computations remain independent of the secret bits, while also
addressing the information leakage present in Trichina's original implementation~\cite{nikovaThresholdImplementationsSideChannel2006}. 

\section{Evaluation Methodology}
Our first milestone will be to design RTL code for a standard BNN MAC without masking, a MAC which masks the multiply operation, 
and a MAC which fully masks the multiply and accumulate operations. Then, our next milestone will be to benchmark the resource utilization
of the hardware designs by compiling each module for synthesis. For this evaluation, we will record the hardware resource usage, energy usage,
and performance of each module. Our final milestone is then to load our module bitstreams onto a an iCE40 FPGA, where they will run in parallel.
We will then use a ChipWhisperer-Husky to run power analysis on each module, where we will perform test vector leakage assessments as well 
as correlation power analysis (CPA) to evaluate the efficacy of masking the weights in each design.

\bibliographystyle{IEEEtran}
\bibliography{references}

\end{document}

